\documentclass{article}
\usepackage[T1]{fontenc}
\usepackage{lmodern}
\usepackage{graphicx}
\usepackage{amsmath,amssymb}
\usepackage[a4paper,margin=2cm]{geometry}
\usepackage[absolute,overlay]{textpos}
\setlength{\TPHorizModule}{1mm}
\setlength{\TPVertModule}{1mm}
\usepackage[indonesian]{babel}
\usepackage{hyperref}
\setlength{\emergencystretch}{2em}
% unit: Halaman 27

\begin{document}

% === Halaman 27 (llm) ===

\textbf{Keterangan:}

\begin{itemize}
    \item $a + b$ adalah penjumlahan dalam modulo $2^w$
    \item $a - b$ adalah pengurangan dalam modulo $2^w$
    \item $a \oplus b$ adalah operasi \textit{bitwise} XOR dari word yang panjangnya $w$ bit
    \item $a \times b$ adalah perkalian dalam modulo $2^w$
\end{itemize}

$a \lll b$ adalah operasi pergeseran $a$ (yang panjangnya $w$ bit) ke kiri secara sirkular sejauh $\log w$ \textit{least significant} bit dari $b$

\vspace{5mm}

$a \ggg b$ adalah operasi pergeseran $a$ (yang panjangnya $w$ bit) ke kanan secara sirkular sejauh $\log w$ \textit{least significant} bit dari $b$

\end{document}
